%! Factoring: GCF and Grouping — Unit 1, Lesson 1.3 — Warm-Up
\documentclass[10pt]{article}
\usepackage{atda-article}
\usepackage{atda-boxes}
\newcommand{\TallMath}[1]{$\displaystyle #1\rule[-1.4em]{0pt}{3.2em}$}
\setlength{\workrowsep}{6pt}   % handwriting room; moves blank and key together

\begin{document}

\pageheader{Unit 1, Lesson 1.3}{Warm-Up --- Running the Distribution Backwards}
% NAMESTRIP: no \namedateperiod here. The student writes their name once, on the
% cover this component is stapled behind. See references/conventions.md.

\vspace{0.15in}

\begin{notesbox}{Four Moves Today's Lesson Runs On}
\small
Item 4 uses your own answer from item 2 --- do them in order. Show your work on all four.

\begin{enumerate}[leftmargin=1.6em, itemsep=6pt, topsep=4pt]

\item \textbf{Spiral --- Lesson 1.2.} Simplify: $\sqrt{48x^3}$
\begin{work}
  \sqrt{48x^3} &= \sqrt{16x^2 \cdot 3x} \\
               &= 4x\sqrt{3x}
\end{work}

\item \textbf{Spiral --- Lesson 1.1.} Multiply out: \quad (a) $5x(3x-7)$ \quad
(b) $(x+2)\left(x^2+3\right)$
\begin{work}
  5x(3x-7) &= 15x^2-35x \\
  (x+2)\left(x^2+3\right) &= x^3+2x^2+3x+6
\end{work}

\item \textbf{Greatest common factor.} Find the largest number that divides both (or all):
\quad $24$ and $36$; \quad $15$ and $25$; \quad $14$, $21$, and $35$
\begin{work}
  24, 36 &\Rightarrow 12 \qquad 15, 25 \Rightarrow 5 \qquad 14, 21, 35 \Rightarrow 7
\end{work}

\item \textbf{Discovery.} Take your answer to 2(b), $x^3+2x^2+3x+6$. Cover the last two terms and
pull out what the \emph{first} two share; then do the same for the last two. What do the two
leftovers have in common?
\begin{work}
  x^3+2x^2+3x+6 &= x^2(x+2)+3(x+2) \\
  \text{Both leftovers are } &(x+2) \text{, so it can be pulled out of the whole thing.}
\end{work}

\end{enumerate}
\end{notesbox}

\end{document}
